\documentclass[aspectratio=169,10pt,fontset=fandol]{ctexbeamer}
\usepackage{amsmath,amssymb,mathtools,bm,mathrsfs}
\usepackage{graphicx,booktabs,array,tabularx,multirow}
\usepackage{tikz,pgfplots,xparse,xstring,mdframed}
\usetikzlibrary{arrows.meta,calc,angles,positioning,shapes.geometric,graphs}
\pgfplotsset{compat=1.18}

\input{workflow_style.tex}
\input{workflow_fonts.tex}
\newenvironment{proofdetail}{\par\noindent\textbf{证明：}\quad}{\par}
\renewenvironment{solution}{\par\noindent\textbf{解：}\quad}{\par}
\title{导数课堂}
\begin{document}
\section{导数及其应用}
% FRAME-ID F001
\typeout{MATH-AGENT-FRAME-BEGIN F001}
\begin{frame}[t,fragile]{目录}
\normalsize
{\color{blue}\normalsize 1. 导数及其应用\par}\vspace{0.22em}
\end{frame}
\typeout{MATH-AGENT-FRAME-END F001}
% FRAME-ID F002
\typeout{MATH-AGENT-FRAME-BEGIN F002}
\begin{frame}[t,fragile]{差商的极限定义}
\normalsize
\begin{block}{定义}
{\normalsize
% SOURCE-BEGIN U0001 0 9
导数是差商的极限。
% SOURCE-END U0001
}
\end{block}
\end{frame}
\typeout{MATH-AGENT-FRAME-END F002}
% FRAME-ID F003
\typeout{MATH-AGENT-FRAME-BEGIN F003}
\begin{frame}[t,fragile]{导数符号与单调性}
\normalsize
\begin{block}{例题}
{\normalsize
% SOURCE-BEGIN U0002 0 23
用导数符号判断 $x^3-3x$ 的单调区间。
% SOURCE-END U0002
}
\end{block}
\end{frame}
\typeout{MATH-AGENT-FRAME-END F003}
\end{document}
